\documentclass[12pt, a4paper]{article} % Classe: article, report o book [8]

% --- Lingua e Codifica (Supporto Italiano) ---
\usepackage[utf8]{inputenc}  % Codifica input
\usepackage[T1]{fontenc}      % Codifica font
\usepackage[italian]{babel}   % Lingua italiana [2]

% --- Impaginazione e Margini ---
\usepackage[margin=2.5cm]{geometry} % Imposta margini (2.5cm su tutti i lati)

% --- Pacchetti Scientifici/Matematici ---
\usepackage{amsmath, amssymb, amsfonts, amsthm} % Matematica avanzata [2]
\usepackage{graphicx} % Per inserire immagini
\usepackage{hyperref} % Per collegamenti ipertestuali
\usepackage{booktabs} % Per tabelle professionali

% --- Formattazione Testo e Altro ---
\usepackage[document]{ragged2e}
\usepackage{xcolor}    % Per usare i colori 

% Titolo e intestazione
\title{Natalia Danielova}
\date{08-01-2026}


\begin{document}
\maketitle

\section{Ruolo}
Natalia Danielova è una commerciale con esperienza di più di 5 anni all'interno
dell'azienda e si occupa di :
\begin{itemize}
      \item Gestione delle richieste di ricambi e assistenza post-vendita.
\end{itemize}

\section{Operazioni comuni}
Le operazioni più comuni svolte da Natalia all'interno dell'applicativo
includono:
\begin{itemize}
      \item Creazione di offerte per ricambi e assistenza (RA).
      \item Creazione degli ordini qualvolta volta un'offerta viene accettata.
      \item Nella creazione delle offerte ricerca e copia di vecchie offerte simili.
      \item Correzione delle traduzioni automatiche fornite da Access.
\end{itemize}

\section{Flusso di lavoro attuale}
Il flusso di lavoro attuale per la gestione dei ricambi e dell'assistenza
post-vendita presenta diverse inefficienze dovute alla frammentazione degli
strumenti e alla mancanza di dati strutturati.

\begin{itemize}
      \item \textbf{Input disorganizzato:} Le richieste dei clienti arrivano principalmente via email in formati non strutturati (testo libero, foto, "collage con Paint"),
            spesso prive del numero di commessa, rendendo difficile l'identificazione dell'impianto e dei componenti.
            Attualmente esiste nel sito web un form per la richiesta ricambi ma non è utilizzato dai clienti abitualmente.
      \item \textbf{Identificazione cliente/impianto:} Spesso il cliente non fornisce il numero di commessa, obbligando l'operatore a cercare manualmente tra i clienti e gli impianti esistenti.
            Una alternativa è richiedre al cliente il numero della commessa impresso sulle etichette affisse nei componenti montati.
      \item \textbf{Template di partenza:} Attualmente non esiste un template di partenza, si parte quindi da vecchi ordini e, per individuare che siano delle assistenze, l'unico discriminante è l'autore.
            In questo caso Natalia in quanto è l'unica che si occupa delle assistenze. Questa pratica permette di portarsi dietri tutte le condizioni di vendita ,
            proprietà del cliente e risparmiare tempo nella compilazione (Attualmente molto oneroso in fatto di tempo).
      \item \textbf{Ricerca ricambi manuale:} L'operatore deve consultare disegni tecnici (Fusion, PDF, cartacei) e database separati (Access, Ad Hoc, Excel , Server "Commesse" ) per identificare i codici dei ricambi.
      \item \textbf{Inserimento del componente:} Una volta individuato il componente , non essendo una corrispondeza univoca dei codici ,
            l'operatore deve trovare la corrispondenza in Access o creare un nuovo codice per il ricambio.
      \item \textbf{Creazione offerta complessa:} Non è facile ne la ricerca ne l'assegnazione di un prezzo se il componente è datato, la compilazione della offerta talvolta richiede calcoli manuali
            (prezzi, pesi, volumi). Spesso si ricorre alla copia di vecchie offerte (RA) o commesse, pratica che trascina errori pregressi.
\end{itemize}

\section{Problemi attuali del sistema}
Dall'intervista a Natalia Danielova sono emerse le seguenti criticità del sistema attuale:
\begin{itemize}
      \item \textbf{Lentezza e complessità:} L'intero processo di gestione delle richieste di ricambi è lento e laborioso, richiede molto tempo per completare un'offerta.
            Questo è dovuto alla necessità di consultare molteplici fonti di dati non integrate. Il programma stesso utilizzato (Access) risulta lento e macchinoso con molti tempi morti.
      \item \textbf{Disallineamento codici:} Esiste una discrepanza tra i codici commerciali e quelli tecnici/gestionali. Questo costringe a ricerche manuali
            su schemi funzionali o database separati per identificare il pezzo corretto (Ad-Hoc , Fusion,  \dots).
      \item \textbf{Mancanza di automazione nei calcoli:}
            \begin{itemize}
                  \item \textbf{Prezzi:} I costi e i prezzi di vendita devono essere spesso calcolati manualmente o verificati su Excel/Ad Hoc, specialmente per componenti su misura o obsoleti.
                  \item \textbf{Logistica:} Non esiste un calcolo automatico di volumi e pesi per il trasporto; l'utente deve ipotizzare i dati basandosi su vecchi DDT o esperienza.
                        Il costo è talvolta calcolato mediante sito di DHL.
                  \item \textbf{Manodopera:} I tempi e i costi per gli interventi tecnici (es. sostituzione saccone) non sono standardizzati e richiedono consultazioni continue con i colleghi (es. Logistica).
            \end{itemize}
      \item \textbf{Componenti elettrici:} La gestione ed inserimento di materiale elettrico è molti difficile per la complessità e varietà di codici esistenti.
            Attualmente si consulta direttamente l'ufficio tecnico elettrico per un elenco dei componenti necessari da offrire.
      \item \textbf{Gestione obsolescenza:} È difficile sapere se un componente venduto anni fa è stato sostituito da una nuova versione o se è obsoleto, rischiando di vendere pezzi non più procurabili.
            Serve un sistema di tracciamento delle sostituzioni e obsolescenze (Es: Fusion).
      \item \textbf{Discrepanza "As Built":} Le modifiche fatte in cantiere (RA di cantiere) o successive alla vendita non sono sempre tracciate in modo unitario, rendendo difficile conoscere la configurazione reale dell'impianto attuale.
            Punto fondamentale per una corretta assistenza senza errori.
      \item \textbf{Sistema di traduzione:} è fondamentale avere un sistema di traduzione integrato per le descrizioni dei componenti e delle offerte,
            attualmente si utilizzano servizi esterni (Google Translate) che richiedono tempo e possono introdurre errori.
            Viene talvolta utilizzato il traduttore automatico di Access.
\end{itemize}

\section{Requisiti Funzionali per il Nuovo Software}

\subsection{Database e integrazione}
Per migliorare l'integrazione dei dati e la gestione delle informazioni, il nuovo sistema deve prevedere:
\begin{itemize}
      \item \textbf{Database unificato:} Il nuovo sistema deve integrare o dialogare con Ad-Hoc (gestionale), Fusion (Pdm) e l'archivio storico per garantire l'univocità del dato.
      \item \textbf{Codifica condivisa:} Allineamento tra codici commerciali e tecnici. Il codice usato per l'offerta deve permettere di risalire automaticamente a costi, distinta base e disegni.
      \item \textbf{Gestione revisioni:} Il sistema deve tracciare le revisioni dei clienti (cambi nome, stabilimento) e dei componenti (sostituzioni, obsolescenza).
      \item \textbf{Codici presenti:} Il sistema deve già avere tutti i codici presenti , la volontà del commerciale è di non doverli inserire manualmente.
            Pensare ad una manutenzione esterna per l'aggiunta e manutenzione di nuovi codici.
      \item \textbf{Integrazione nel flusso dell'ordine:} Attualmente esistono 3 codici fodnamentali per la gestione dell'ordine :
            \begin{itemize}
                  \item \textbf{Codice offerta (RA):} Codice univoco dell'offerta commerciale.
                  \item \textbf{Codice offerta (RA):} Codice univoco dell'offerta trasformata in ordine.
                  \item \textbf{Codice distinta (RA):} Codice univoco della distinta base gestionale (Ad-Hoc) dell'ordine.
            \end{itemize}
            Il nuovo sistema deve integrare questi 3 codici in modo chiaro e visibile per l'operatore.
            Facilitando le ricerche e la tracciabilità dell'ordine.
\end{itemize}

\subsection{Ricerca e navigazione ricambi}
Per facilitare la ricerca e la selezione dei ricambi, il sistema deve prevedere:
\begin{itemize}
      \item \textbf{Ricerca multicriterio:} Necessità di cercare per Numero di Commessa (fondamentale), Cliente, Stabilimento, e Tag Tecnici (es. "valvola a farfalla").
      \item \textbf{Struttura ad albero/gerarchica:} Possibilità di esplorare la distinta base (BOM) scendendo nei livelli (esplosione componenti) per individuare il singolo ricambio.
      \item \textbf{Filtri per "Ricambi tipici":} Implementazione di filtri che mostrino, per una data famiglia di macchine (es. Lift, Vibrovaglio), solo i componenti tipicamente venduti come ricambi,
            i codici devono essere univoci con i codici tecnici rendendo dopo l'idividuazione del componente veloce l'assistenza.
\end{itemize}

\subsection{Creazione e gestione offerte (RA)}
Per semplificare e velocizzare la creazione delle offerte di ricambi e assistenza, il sistema deve prevedere:
\begin{itemize}
      \item \textbf{Automazione descrizioni e lingue:} Compilazione automatica delle descrizioni in lingua (traduzioni pre-caricate) senza dover usare traduttori esterni.
      \item \textbf{Calcolo prezzi automatizzato:} Il sistema deve pescare il costo industriale da Ad-Hoc e applicare le maggiorazioni o i listini specifici per cliente/settore.
      \item \textbf{Numero univoco offerta:} Mantenimento di un numero univoco che segue la pratica dalla richiesta iniziale all'ordine confermato, per evitare duplicazioni e facilitare la rintracciabilità.
            Aggiunta di uno stato di avanzamento \textit{(es. "In Revisione", "Inviata", "Accettata", "Rifiutata")} facilitando l'interpretazione.
      \item \textbf{Distinzione tipologie RA:} Gestione differenziata ma tracciabile per "Ricambi", "Ampliamenti" e "RA di Cantiere". Al momento le assistenze prima citate sono sotto un unico cappello del "VSS".
\end{itemize}

\subsection{Funzionalità avanzate e UX}
Per migliorare l'esperienza utente e l'efficienza operativa, il sistema deve prevedere:
\begin{itemize}
      \item \textbf{Schermata offerta:} Importante riporatare tutte le infomazioni attuali del cliente (Indirizzo , Nome , Numero di commessa associata) in modo chiaro e visibile.
      \item \textbf{Scheramata dedicata allo storico:} Una sezione dove poter visualizzare tutte le offerte/ordini precedenti associati al cliente o alla commessa selezionata.
      \item \textbf{Schermata offerta intuitiva:} Interfaccia user-friendly che guidi l'utente passo-passo nella compilazione dell'offerta, minimizzando errori e omissioni.
      \item \textbf{Visualizzazione dati completa:} Interfaccia che mostri tutte le informazioni rilevanti (es. dettagli tecnici completi come "filtrante/non filtrante") senza dover scorrere o aprire popup.
\end{itemize}

\subsection{Integrazioni future e scalabilità}
Per garantire la scalabilità e l'integrazione con altri sistemi futuri, il sistema deve prevedere:
\begin{itemize}
      \item \textbf{Integrazione disegni/foto (Configuratore):} In prospettiva futura, un sistema visuale dove il cliente o l'operatore possa cliccare sul componente nel disegno 3D/2D per identificarlo univocamente,
            riducendo l'uso di foto e interpretazioni manuali.
      \item \textbf{Integrazione con gestionale produzione:} Il distema non è integrato con TWproject (gestionale produzione) ,
            ma sarebbe utile avere un collegamento per verificare lo stato di avanzamento delle commesse.
      \item \textbf{Integrazione Fusion:} Il sistema sarebbe utile essere integrato con Fusion per la visualizzazione dei disegni tecnici e la gestione delle modifiche.
\end{itemize}

\section{Numeri utili}
Per capire gli attuali tempi di lavoro e il volume di offerte gestite da Natalia, si riportano i seguenti dati:
\begin{itemize}
      \item \textbf{Offerte VS Ordini} Attualmente natalia ha creato circa 57000 offerte (RA) di cui circa 41000 si sono trasformate in ordini effettivi.
      \item \textbf{Tempo base per una semplice offerta:} Attualmente per una commessa tipo dove si conosce il numero di commessa e i codici da inserire
            il tempo medio per la creazione di una offerta è di circa 15 minuti.
            Il tempo sale notevolmente se sono da ricercare le informazioni mancanti.
\end{itemize}
\end{document}